\section{Results and Discussion}
\label{sec:results}

% =============================================================================

% -----------------------------------------------------------------------------
% PER-SECTION PATTERN:
%   1. one paragraph: what is being compared and how
%   2. table (one per subsection in most cases)
%   3. one or two paragraphs: what we observed and why we picked the winner
%   4. last sentence: name the locked-in choice, since the next section assumes it
%
% NUMBERS in the tables below come from tuning_notes.txt and are pre-filled so
% you can sanity-check rather than transcribe. Verify against the latest runs
% before submission.
%
% TABLE BUDGET: aim for 6-8 tables total in this chapter. Most subsections have
% one table; 5.6 has two. If space gets tight, the BM25 grid (5.2) and the
% alpha sweep (5.6.1) can collapse into prose.
%
% MULTILINGUAL: report per-language EN / FR / DE columns plus Avg. Do NOT use
% per-month line plots (those were for LongEval, which is temporal). Our axis
% is language, not time.
% =============================================================================

In this section we report the experiments that led to the final lexical and hybrid configurations of the system. Because the search space is much larger than what we could grid-search, we tuned the components in order, fixing the winner of each phase before moving on:

\begin{enumerate}
    \item Section~\ref{subsec:res-analyzer} fixes the analyzer pipeline (tokenizer, stemmer, stop list, filters).
    \item Section~\ref{subsec:res-model} fixes the retrieval model and BM25 parameters.
    \item Section~\ref{subsec:res-fielded} fixes the fielded query (boosts, query mode, IDF pruning).
    \item Section~\ref{subsec:res-expansion} reports two expansion strategies (WordNet, PRF) that did not work.
    \item Section~\ref{subsec:res-ablation} ablates the components of the final lexical system.
    \item Section~\ref{subsec:res-hybrid} reports the dense fusion and reranking results.
    \item Section~\ref{subsec:res-final} compares the runs we submitted.
\end{enumerate}

All scores are MRR@5 on the dev split, reported per language (EN, FR, DE) with the average across languages. Hybrid and reranking experiments operate on the top-50 candidates of the lexical run. The dataset, dev splits, hardware, and evaluation tooling are described in Section~\ref{sec:setup}.

% TODO (anyone): if there are non-trivial caveats about dev data, evaluation tooling,
% or qrels processing, raise them here in one short paragraph. Skip if there are none.

\subsection{Lexical Analyzer Tuning}
\label{subsec:res-analyzer}

% =============================================================================
% SOURCE: tuning_notes.txt §3 (FR/DE artefacts), §4 (stop-list refinement),
%         §5 (tokenizer), §6 (stemmer), §7 (shingle), §8 (English possessive)
% GOAL:   establish the analyzer config that all later sections build on.
% TABLE:  one combined comparison (tokenizer x stemmer + stop-list ablation).
% COVER:
%   - StandardTokenizer beats WhitespaceTokenizer by ~5pp average
%   - the stop-list contribution (~2pp) is larger than the gap between stemmers (<0.3pp)
%   - KStem narrowly best on average; preferred for being a conservative stemmer
%     ("vaccination"/"vaccinate" conflate, "corona"/"coronary" do not)
%   - shingles hurt 3-5pp -> off (one paragraph, no table needed; explain that
%     BM25 distributes IDF across unigram + bigram, claims already carry phrase
%     signal so shingles dilute)
%   - EnglishPossessiveFilter neutral, kept on (one sentence)
% AVOID:
%   - re-stating the analyzer architecture (already in methodology)
%   - the full §3 vocabulary cleanup story; one passing sentence about FR/DE
%     synonym/POS filters being disabled is enough
% =============================================================================

We start by fixing the analyzer pipeline that the rest of the chapter builds on: the tokenizer, the stop list, the stemming strategy, and a few token filters.

% TODO: expand the opener above into a full paragraph naming what is being compared.

\begin{table}[!ht]
    \caption{Tokenizer and stemmer comparison on the EN/FR/DE dev sets, MRR@5. Variants share the standard pipeline (lowercase, ICU folding, English stop list) and differ only in the stemmer or tokenizer.}
    \label{tab:stemmer-comparison}
    \centering
    \begin{tabular}{lcccc}
        \toprule
        Configuration & EN & FR & DE & Avg \\
        \midrule
        Standard, no stop, no stem        & 0.4826 & 0.4453 & 0.3014 & 0.4098 \\
        Standard, stop, no stem           & 0.4864 & 0.4674 & 0.3225 & 0.4254 \\
        Standard, stop, EnglishMinimal    & 0.4894 & 0.4817 & 0.3304 & 0.4338 \\
        Standard, stop, Snowball          & 0.4869 & 0.4785 & 0.3358 & 0.4337 \\
        Standard, stop, Porter            & 0.4867 & 0.4801 & 0.3381 & 0.4350 \\
        \textbf{Standard, stop, KStem}    & \textbf{0.4883} & \textbf{0.4901} & \textbf{0.3354} & \textbf{0.4379} \\
        \midrule
        Whitespace, stop, KStem           & 0.4422 & 0.4087 & 0.2746 & 0.3752 \\
        \bottomrule
    \end{tabular}
\end{table}

% TODO: 1-2 paragraphs interpreting the table. Cover (in order):
%   - tokenizer: Standard beats Whitespace by ~5pp because Whitespace leaves
%     punctuation attached and fails to split hyphenated biomedical compounds.
%   - stop list contribution (~2pp on FR) is larger than any stemmer-vs-stemmer gap.
%   - all four stemmers are within 0.3pp on EN; KStem narrowly best on average,
%     and chosen because it is the most conservative (avoids over-stemming).
%   - last sentence: state the locked-in pipeline used in 5.2 onwards.

% TODO: one short paragraph on shingles (§7).
% Key numbers: 2-shingle = 0.3855 avg, 3-shingle = 0.4040 avg, vs baseline 0.4380.
% Reason: BM25 spreads IDF across unigrams and bigrams; claims are already
% complete sentences so shingles add redundant tokens that dilute precise lexical
% matches. Decision: shingles off.

% TODO: one line on EnglishPossessiveFilter (§8). Difference < 0.01pp; kept on
% since it has no overhead and is conceptually correct.

\subsection{Retrieval Model and BM25 Parameters}
\label{subsec:res-model}

% =============================================================================
% SOURCE: tuning_notes.txt §11 (similarity comparison), §12 (k1 x b grid)
% GOAL:   pick BM25 over alternatives, fix k1 and b.
% TABLES: one model comparison table (below). The k1 x b grid is described in
%         prose; if you have space, add a second table or heatmap, but the
%         column at b=0.75 is essentially flat across k1 in [0.8, 1.5] so prose
%         is enough.
% COVER:
%   - BM25 wins by ~1.2pp over best LM alternative
%   - Classic TF-IDF worst (no length normalisation -> long abstracts dominate)
%   - LMD degrades as mu grows (scientific abstracts dense in specific vocab)
%   - in the BM25 grid, b dominates k1 entirely:
%       moving b from 0 to 0.75 adds ~5pp at every k1
%       at b=0.75 the entire k1 column varies by < 0.001
%   - Lucene defaults (k1=1.2, b=0.75) are within 0.0015 of optimum;
%     winner is k1=1.0, b=0.75 with avg MRR@5 = 0.4357
% AVOID:
%   - the full 7x5 BM25 heatmap (noisy and the column is flat); summarise instead
% =============================================================================

With the analyzer fixed, we choose a similarity function and tune its parameters. We compare BM25 against Lucene's other built-in similarity models, then sweep BM25's $k_1$ and $b$ over a small grid.

% TODO: expand the opener above; describe what is being compared.

\begin{table}[!ht]
    \caption{Retrieval model comparison, MRR@5 on the dev sets. BM25 uses Lucene defaults ($k_1{=}1.2$, $b{=}0.75$); only the most informative LM Dirichlet and LM Jelinek-Mercer rows are shown.}
    \label{tab:model-comparison}
    \centering
    \begin{tabular}{lcccc}
        \toprule
        Model & EN & FR & DE & Avg \\
        \midrule
        \textbf{BM25} ($k_1{=}1.2$, $b{=}0.75$)     & \textbf{0.4838} & \textbf{0.4887} & \textbf{0.3342} & \textbf{0.4356} \\
        Classic (TF-IDF)                             & 0.4130 & 0.3948 & 0.2715 & 0.3597 \\
        LM Dirichlet ($\mu{=}500$)                   & 0.4683 & 0.4652 & 0.3357 & 0.4231 \\
        LM Dirichlet ($\mu{=}1000$)                  & 0.4526 & 0.4434 & 0.3098 & 0.4019 \\
        LM Jelinek-Mercer ($\lambda{=}0.5$)          & 0.4784 & 0.4706 & 0.3220 & 0.4236 \\
        \bottomrule
    \end{tabular}
\end{table}

% TODO: paragraph interpreting the model comparison. Order: BM25 winner ->
% Classic worst (no length norm) -> LMD trends with mu -> LMJM never matches BM25.

% TODO: paragraph on the BM25 k1 x b sweep. Communicate:
%   - swept k1 in {0.5..3.0} and b in {0.0..1.0}
%   - winner k1=1.0, b=0.75 (avg MRR@5 = 0.4357)
%   - b dominates k1 entirely; at b=0.75 the k1 column varies by < 0.001
%   - Lucene defaults are within 0.0015 of optimum; the change is +0.15pp
% End with the locked-in BM25 settings used for the rest of the chapter.

\subsection{Fielded Search and Query Construction}
\label{subsec:res-fielded}

% =============================================================================
% SOURCE: tuning_notes.txt §13 (title/abstract boost), §14 (authors/venue),
%         §15 (must vs should), §16 (topKQueryTerms IDF pruning)
% GOAL:   fix the structured query: which fields to boost, how queries are assembled.
% TABLES: one boost-sweep table covering title and abstract; one topK pruning table.
%         queryMode (§15) is one paragraph, no table - the must-mode collapse is
%         too dramatic to need a table.
% ORDER:  title/abstract -> authors/venue -> queryMode -> topK
% COVER:
%   - title^2 abstract^3 wins (the ratio matters more than absolute values).
%     Claims are tweet-style paraphrases of findings, so abstracts carry more
%     signal than titles.
%   - authors^1.5 helps (~+1.5pp); tweets sometimes mention researcher names.
%   - venue boost is neutral (~0pp); not adopted.
%   - queryMode=must collapses to ~0.07 avg; even must=2 drops below 0.10.
%     Claims never quote paper text verbatim, so any hard term constraint
%     destroys recall. SHOULD is the only viable mode.
%   - topK=30 plateau: 30 = 40 = 50 = peak. At topK=10-15 genuine scientific
%     terms are dropped.
% AVOID:
%   - re-explaining the boost mechanism (already in methodology)
% =============================================================================

We now tune how the query is assembled: which fields to boost, whether terms are required or optional, and how aggressively to prune the query.

% TODO: expand the opener above into a full paragraph.

\begin{table}[!ht]
    \caption{Field boost sweeps on the EN/FR/DE dev sets, MRR@5. Part A varies the title boost with abstract fixed at 1.0; Part B varies the abstract boost with title fixed at 2.0; Part C adds the authors boost on top of title$^2$ abstract$^3$.}
    \label{tab:field-boosts}
    \centering
    \begin{tabular}{lcccc}
        \toprule
        Configuration & EN & FR & DE & Avg \\
        \midrule
        \multicolumn{5}{l}{\textit{Part A: title boost (abstract$^1$)}} \\
        title$^{1.0}$              & 0.4997 & 0.5026 & 0.3577 & 0.4533 \\
        title$^{1.5}$              & 0.4929 & 0.4964 & 0.3475 & 0.4456 \\
        title$^{2.0}$              & 0.4838 & 0.4882 & 0.3352 & 0.4357 \\
        title$^{3.0}$              & 0.4635 & 0.4567 & 0.3159 & 0.4121 \\
        \midrule
        \multicolumn{5}{l}{\textit{Part B: abstract boost (title$^2$)}} \\
        abstract$^{1.0}$           & 0.4838 & 0.4882 & 0.3352 & 0.4357 \\
        abstract$^{1.5}$           & 0.4956 & 0.5000 & 0.3506 & 0.4487 \\
        abstract$^{2.0}$           & 0.4997 & 0.5026 & 0.3577 & 0.4533 \\
        abstract$^{3.0}$           & 0.4990 & 0.5066 & 0.3604 & 0.4553 \\
        \midrule
        \multicolumn{5}{l}{\textit{Part C: authors boost (title$^2$ abstract$^3$, topK${=}30$)}} \\
        authors$^{0.5}$            & 0.4888 & 0.4938 & 0.3373 & 0.4400 \\
        authors$^{1.0}$            & 0.4990 & 0.5066 & 0.3604 & 0.4693 \\
        \textbf{authors$^{1.5}$}   & \textbf{0.5172} & \textbf{0.5145} & \textbf{0.3841} & \textbf{0.4732} \\
        authors$^{2.0}$            & 0.5160 & 0.5130 & 0.3820 & 0.4715 \\
        authors$^{3.0}$            & 0.5050 & 0.5020 & 0.3754 & 0.4608 \\
        \bottomrule
    \end{tabular}
\end{table}

% TODO: 2 paragraphs.
%   1. Title and abstract boosts: title^2 abstract^3 (ratio 2:3) is best. Explain
%      WHY: claims paraphrase findings, not titles, so the abstract is the primary
%      evidence field. Note that abstract-only (title->0) still scores 0.4061, so
%      titles do contribute, just less.
%   2. Authors and venue: authors^1.5 peaks (DE benefits most - tweets translated
%      from German more often name researchers). Higher than 1.5 hurts as
%      mismatched proper nouns dominate. Venue swept to ~0.5; near-zero effect,
%      not adopted.

% TODO: short paragraph (3-4 lines) on queryMode (§15). Three lines is enough:
%   - tested SHOULD (default), mixed (must=N highest-IDF + SHOULD rest), MUST (all)
%   - mixed must=2 drops to avg 0.087; full MUST collapses to 0.007
%   - claims never quote paper text verbatim, so any hard constraint destroys recall
%   - SHOULD is the only viable mode; used everywhere else in this chapter

\begin{table}[!ht]
    \caption{IDF-based query pruning on top of the fielded query, MRR@5. Each claim is tokenised with the English analyzer; only the top-$K$ tokens by IDF are passed as the query.}
    \label{tab:topk-pruning}
    \centering
    \begin{tabular}{lcccc}
        \toprule
        $K$ & EN & FR & DE & Avg \\
        \midrule
        all (no pruning)   & 0.4838 & 0.4882 & 0.3352 & 0.4357 \\
        10                 & 0.4115 & 0.4189 & 0.2887 & 0.3730 \\
        15                 & 0.4696 & 0.4658 & 0.3182 & 0.4179 \\
        20                 & 0.4914 & 0.4798 & 0.3379 & 0.4364 \\
        \textbf{30}        & \textbf{0.4956} & \textbf{0.4860} & \textbf{0.3395} & \textbf{0.4404} \\
        40                 & 0.4956 & 0.4860 & 0.3395 & 0.4404 \\
        50                 & 0.4956 & 0.4860 & 0.3395 & 0.4404 \\
        \bottomrule
    \end{tabular}
\end{table}

% TODO: short paragraph on topK. Plateau at 30; smaller values drop genuine
% scientific terms. Most claims contain fewer than 30 distinct content tokens
% after stop-word filtering, hence the flat tail.

\subsection{Failed Expansion Approaches}
\label{subsec:res-expansion}

% =============================================================================
% SOURCE: tuning_notes.txt §9 (WordNet), §10 (PRF). Methodology already
%         describes both mechanisms.
% GOAL:   report negative results and explain WHY they failed (this is the
%         interesting part: both fail for the same root cause).
% TABLE:  one combined table with baseline + WordNet variants + PRF variants.
% COVER:
%   - both hurt monotonically as expansion grows
%   - common cause: claims are already long, specific scientific sentences;
%     adding terms dilutes BM25's IDF-weighted match on the precise originals
%   - WordNet additionally suffers from domain mismatch (general-vocabulary lexicon)
%   - PRF additionally doubles retrieval time and topic-drifts toward adjacent papers
%   - CONTRAST with §5.6: BGE-M3 fusion DOES help. The difference is that
%     embeddings capture semantic similarity directly, whereas expansion adds raw
%     tokens that interact badly with BM25 length normalisation.
%     This contrast is the single most useful insight in the section. Make it
%     explicit in the closing paragraph.
% =============================================================================

We tested two query expansion strategies that the methodology supports, WordNet synonyms and pseudo-relevance feedback. Both hurt retrieval, and for the same underlying reason.

% TODO: expand the opener above; explain why we tried these and what the methodology promised.

\begin{table}[!ht]
    \caption{Query expansion experiments, MRR@5. All runs use the lexical configuration fixed in Sections \ref{subsec:res-analyzer}--\ref{subsec:res-fielded}.}
    \label{tab:expansion}
    \centering
    \begin{tabular}{lcccc}
        \toprule
        Configuration & EN & FR & DE & Avg \\
        \midrule
        Lexical baseline (no expansion)         & 0.4883 & 0.4901 & 0.3354 & 0.4379 \\
        \midrule
        WordNet, max 1 syn/word                  & 0.4582 & 0.4352 & 0.3040 & 0.3991 \\
        WordNet, max 3 syn/word                  & 0.4262 & 0.3993 & 0.2829 & 0.3695 \\
        WordNet, max 5 syn/word                  & 0.4047 & 0.3759 & 0.2591 & 0.3466 \\
        \midrule
        PRF (3 docs, 3 terms)                    & 0.4825 & 0.4763 & 0.3351 & 0.4313 \\
        PRF (5 docs, 3 terms)                    & 0.4829 & 0.4790 & 0.3310 & 0.4310 \\
        PRF (10 docs, 10 terms)                  & 0.4776 & 0.4692 & 0.3317 & 0.4262 \\
        \bottomrule
    \end{tabular}
\end{table}

% TODO: 2 paragraphs.
%   1. WordNet: monotonic degradation as more synonyms are added. Three causes:
%      domain mismatch (scientific terms absent from WordNet), specificity loss
%      (synonyms of "system" point to general everyday concepts), and BM25
%      dilution (extra tokens lower the IDF contribution of the originals).
%   2. PRF: same dilution mechanism plus a 2x retrieval cost. The top retrieved
%      docs contain related-but-different topics in the same domain (e.g. nearby
%      COVID papers), so feeding those terms back drifts the query toward a
%      broader topic cluster.
% Closing: both disabled in the submitted runs. Foreshadow §5.6: dense fusion
% addresses the same vocabulary-mismatch problem without the dilution.

\subsection{Ablation of the Lexical System}
\label{subsec:res-ablation}

% =============================================================================
% SOURCE: tuning_notes.txt §17. This is the centrepiece table - keep it whole.
% GOAL:   show the marginal contribution of each lexical component on top of
%         the full best lexical configuration.
% TABLE:  one table, one row per component removed, with delta column.
% COVER:
%   - abstractBoost (abs^1 -> abs^3) is the strongest single contributor (-0.025)
%   - authorsBoost (-0.015)
%   - topKQueryTerms (-0.003): small but positive, kept for compactness
%   - venue near zero, not adopted
%   - components slightly sub-additive (sum 0.043, total 0.035): expected for
%     non-orthogonal retrieval improvements
%   - DE benefits disproportionately from authorsBoost (-0.020 vs -0.009 EN)
%   - queryMode=must rows confirm §5.3: catastrophic in every language
% AVOID:
%   - re-explaining what each component does (already in 5.3)
% =============================================================================

To isolate the contribution of each lexical component, we ablate them one at a time on top of the full best configuration, sharing a single index across runs.

% TODO: expand the opener above into a paragraph describing what the table shows and how it was constructed.

\begin{table}[!ht]
    \caption{Ablation of the lexical system, MRR@5. The reference is the full configuration (BM25 $k_1{=}1.0$, $b{=}0.75$, title$^2$ abstract$^3$ authors$^{1.5}$, topK${=}30$). Each row removes or changes exactly one component.}
    \label{tab:ablation}
    \centering
    \begin{tabular}{lccccr}
        \toprule
        Configuration & EN & FR & DE & Avg & $\Delta$ Avg \\
        \midrule
        Full best (reference)                    & 0.5172 & 0.5145 & 0.3841 & 0.4719 & --- \\
        \midrule
        $-$ authorsBoost (set to 0)              & 0.5086 & 0.4973 & 0.3646 & 0.4568 & $-0.0151$ \\
        $-$ topKQueryTerms (use all terms)       & 0.5100 & 0.5131 & 0.3826 & 0.4686 & $-0.0034$ \\
        $-$ abstractBoost (restore abs$^1$)      & 0.4983 & 0.4885 & 0.3548 & 0.4472 & $-0.0247$ \\
        $-$ all three (simple baseline)          & 0.4874 & 0.4919 & 0.3323 & 0.4372 & $-0.0347$ \\
        \midrule
        $+$ venueBoost ${=}\,0.5$                & 0.5174 & 0.5155 & 0.3835 & 0.4721 & $+0.0002$ \\
        queryMode = mixed (must${=}\,2$)         & 0.0752 & 0.1233 & 0.0635 & 0.0873 & $-0.3846$ \\
        queryMode = must (all)                   & 0.0151 & 0.0043 & 0.0026 & 0.0073 & $-0.4646$ \\
        \bottomrule
    \end{tabular}
\end{table}

% TODO: 2 paragraphs.
%   1. Ranked component contributions: abstractBoost > authorsBoost > topK >> venue.
%      Note the DE-specific authors effect (-0.0195 vs -0.0086 EN).
%      Note that the components are slightly sub-additive (0.043 sum vs 0.035
%      observed total) - expected when the components are not orthogonal.
%   2. One short paragraph on queryMode rows: catastrophic in every language,
%      confirms the choice of SHOULD made in §5.3. No claim text appears verbatim
%      in any abstract, so no hard term constraint is safe.

\subsection{Hybrid Retrieval: Lexical and Dense Fusion}
\label{subsec:res-hybrid}

% =============================================================================
% SOURCE: tuning_notes.txt §19 (BGE-M3 fusion + alpha sweep), §20 (SPECTER2),
%         §21 (cross-space; exploratory).
%
% IMPORTANT GAP: project memory mentions a MiniLM cross-encoder reranker on top
% of fusion (rerankAlpha=0.6, top 50). This is NOT in the tuning notes. Whoever
% writes this section needs to pull MiniLM dev-set numbers from the latest runs
% and add subsubsection 5.6.3. Without those numbers, the section ends at the
% SPECTER2 comparison.
%
% GOAL:   show that BGE-M3 fusion adds substantial gain on top of the lexical
%         system (~+5.7pp avg) and that SPECTER2 underperforms BGE-M3 here.
% TABLES: two tables - alpha sweep summary + SPECTER2 vs BGE-M3.
% COVER:
%   - fusion >> reranking-only at the optimum alpha (+5.7pp vs +1.2pp)
%   - alpha curve flat between 0.67 and 0.75; alpha=0.70 is the practical choice
%   - alpha=0.70 means BGE-M3 weighted 2.3x more than BM25 - BM25 near its ceiling
%   - SPECTER2 (domain-specific) loses to BGE-M3 (multilingual general) by ~2.4pp
%   - this contrasts with the §5.4 expansion failures: embeddings capture
%     semantic similarity directly, which is the right tool for tweet -> abstract
%   - cross-space PCA experiment (§21) shows SPECTER2 doc representations are
%     also not better - keep this to ONE paragraph, do NOT give it a subsubsection
% STRUCTURE:
%   5.6.1 BGE-M3 reranking vs fusion (table + alpha sweep summary)
%   5.6.2 SPECTER2 vs BGE-M3 (table + paragraph)
%   5.6.3 Cross-encoder reranking on the fused list   <-- only if MiniLM data added
% =============================================================================

On top of the lexical system fixed in the previous sections, we add a dense retrieval signal from a multilingual sentence encoder and combine the two scores. We compare two encoder families and two ways of combining the scores.

\subsubsection{BGE-M3 Reranking and Fusion}

We first compare two ways of using BGE-M3 embeddings: pure reranking, which reorders the top-50 lexical hits by cosine similarity, and score fusion, which weights the BM25 and cosine scores against each other within each query.

% TODO: expand the opener above; cover the setup (top-50 candidates from the
% lexical run), what reranking-only and fusion mean, the fusion formula, alpha
% definition. The fusion formula: combined = (1 - alpha) * BM25_norm + alpha *
% cosine, with min-max normalisation within each query's top-50 before combining.

\begin{table}[!ht]
    \caption{Lexical baseline vs.\ BGE-M3 reranking and fusion, MRR@5. Reranking reorders the top 50 lexical hits by cosine; fusion combines $(1{-}\alpha)\cdot \text{BM25}_{\text{norm}} + \alpha\cdot \text{cosine}$ with both scores min-max normalised within each query's top-50.}
    \label{tab:bge-fusion}
    \centering
    \begin{tabular}{lcccc}
        \toprule
        Method & EN & FR & DE & Avg \\
        \midrule
        Lexical baseline                          & 0.5172 & 0.5145 & 0.3841 & 0.4719 \\
        BGE-M3 reranking only                     & 0.5256 & 0.5280 & 0.3967 & 0.4834 \\
        \textbf{BGE-M3 fusion} ($\alpha{=}0.70$)  & \textbf{0.5768} & \textbf{0.5770} & \textbf{0.4331} & \textbf{0.5290} \\
        \bottomrule
    \end{tabular}
\end{table}

% TODO: paragraph on the alpha sweep. Either include a small sweep table or
% describe in prose. Key facts:
%   - per-language peaks all in [0.69, 0.73]: EN ~ 0.73, FR ~ 0.73, DE ~ 0.69
%   - the curve is flat from alpha=0.67 to alpha=0.75 (within < 0.002 of peak)
%   - alpha=0.70 is within 0.002 of every per-language peak; chosen as the
%     single practical setting

% TODO: 1-2 paragraphs interpreting:
%   - reranking-only is weak (+1.2pp) because the top-50 already contains the
%     correct doc near the top in many cases - little re-ordering headroom
%   - fusion (+5.7pp) recovers candidates ranked 6-50 that BM25 underscored but
%     BGE-M3 correctly identifies as semantically close
%   - alpha=0.70 weights the embedding 2.3x BM25, suggesting BM25 is near its
%     ceiling on this task and the embedding captures complementary evidence

\subsubsection{SPECTER2 vs BGE-M3}

BGE-M3 is a general-purpose multilingual encoder. We also tested SPECTER2, a domain-specific encoder trained on scientific paper-to-paper citation proximity, to see whether the scientific specialisation outweighs the loss of multilingual coverage. It does not.

% TODO: expand the opener above; the hypothesis was that the scientific encoder
% would beat the general multilingual one on scientific retrieval, but the
% result was the opposite.

\begin{table}[!ht]
    \caption{SPECTER2 vs.\ BGE-M3 at their respective best fusion weights. Both use the same 50-candidate input lists from the lexical run.}
    \label{tab:specter-vs-bge}
    \centering
    \begin{tabular}{lcccc}
        \toprule
        Method & EN & FR & DE & Avg \\
        \midrule
        Lexical baseline                           & 0.5172 & 0.5145 & 0.3841 & 0.4719 \\
        SPECTER2 reranking only                    & 0.4451 & 0.4233 & 0.3377 & 0.4020 \\
        SPECTER2 fusion ($\alpha{=}0.60$)          & 0.5524 & 0.5475 & 0.4150 & 0.5050 \\
        \textbf{BGE-M3 fusion} ($\alpha{=}0.70$)   & \textbf{0.5768} & \textbf{0.5770} & \textbf{0.4331} & \textbf{0.5290} \\
        \bottomrule
    \end{tabular}
\end{table}

% TODO: paragraph explaining why BGE-M3 wins:
%   - SPECTER2 is trained on paper-to-paper citation proximity; its query-side
%     encoder (the adhoc_query adapter) has never seen informal social media text.
%   - BGE-M3 is trained on broad text pairs including informal/formal cross-domain
%     matching - directly bridges tweet paraphrase -> formal abstract.
%   - SPECTER2 needs a lower alpha (0.60), meaning BM25 is weighted more: the
%     SPECTER2 signal is weaker, not stronger.

% TODO: ONE paragraph on the cross-space experiment (§21). Do NOT give it a
% subsubsection. Cover:
%   - we paired SPECTER2 doc embeddings with BGE-M3 query embeddings (PCA-aligned
%     to a common 768-d space)
%   - cross-space cosine collapses to ~0.025 MRR@5 (PCA into a foreign space is
%     geometrically meaningless; this run is purely diagnostic)
%   - best fusion alpha drops to 0.1 and the fused result is indistinguishable
%     from the lexical baseline (+0.0001 to +0.01)
% Conclusion: the failure is not just SPECTER2's query side; the document
% representations are also not aligned with what tweet-style claim retrieval needs.
% BGE-M3 at alpha=0.70 is the encoder of choice; the embedding model comparison
% is closed.

% --- IF MiniLM RERANKER NUMBERS ARE AVAILABLE ---
% \subsubsection{Cross-Encoder Reranking}
%
% TODO: only include if the submitted system uses MiniLM on top of fusion.
% Per project notes: best config = hybrid + rerank, top 50, rerankAlpha=0.6.
% Numbers needed: dev MRR@5 for the hybrid system with and without the reranker,
% per language. Pull from the latest runs - NOT in tuning_notes.txt.

\subsection{Final Submitted Runs}
\label{subsec:res-final}

% =============================================================================
% SOURCE: pulls together winners from §5.1-5.6.
% GOAL:   present the runs we submit and their dev-set scores in one place.
% TABLE:  one summary table.
% COVER:
%   - which runs we submit and why
%   - quote per-language scores so multilingual robustness is visible
%   - flag whether MiniLM rerank is part of the primary submission
% OPTIONAL: per-language bar chart comparing 3-4 submitted runs (analogous to
%           3DS2A's Figure 12). Only add if it shows something the table doesn't.
% =============================================================================

Combining the winners of each phase above gives the configurations we submit to the official evaluation.

% TODO: expand the opener above; state which runs were submitted and why.
% Likely candidates (confirm with the team):
%   1. Lexical only (best lexical config) - cheap baseline
%   2. Hybrid (lexical + BGE-M3 fusion at alpha=0.70) - main submission
%   3. Hybrid + MiniLM rerank - if numbers available, primary submission

\begin{table}[!ht]
    \caption{Submitted runs and dev-set MRR@5 per language.}
    \label{tab:final-runs}
    \centering
    \begin{tabular}{llcccc}
        \toprule
        Run & Description & EN & FR & DE & Avg \\
        \midrule
        Run 1 & Lexical only (best lexical config)              & 0.5172 & 0.5145 & 0.3841 & 0.4719 \\
        Run 2 & Hybrid: lexical $+$ BGE-M3 fusion ($\alpha{=}0.70$) & 0.5768 & 0.5770 & 0.4331 & 0.5290 \\
        % Run 3 & Hybrid $+$ MiniLM rerank ($\alpha_{\text{rerank}}{=}0.6$) & ? & ? & ? & ? \\  % TODO if applicable
        \bottomrule
    \end{tabular}
\end{table}

% TODO: short closing paragraph: which run is the primary submission and why,
% any caveats on test-set extrapolation, transition into the conclusion.
